\subsection{Research Problem, Aim, and Research Questions}

\guidance{When you have established the background in Section 1.1, you switch over to this subsection that should explain the specific problem statement and/or knowledge gap that you have identified. Then, highlight the aim of this project, in terms of what part of the research problem/gap you want to address. Finally, formulate some specific questions that are aligned with the aims, that have uncertain answers, and that you will seek answers to.}

\guidance{These questions should be formulated with sufficient precision and narrowness that you can find answers to them within the scope of the project. A bad example is “What is the best way to achieve X?” since it neither specifies how to measure what is best nor limits the number of ways that must be analyzed. Consider your research questions from the SMART criteria: Are they Specific, Measurable, Achievable, Relevant, and Time-bound? If a research question can be stated to have a true or false answer, then you can also provide your hypothesis of what the answer might be.}

\placeholder{[Problem statement / research gap, research aims, research questions, and optional hypotheses (statement or statements to examine during the project)]}